\paragraph{The forms in the selected realization.}\label{frame:forms}
At the distinguished embedding take the reduced source to be the
rational preimage of the range of the adjoint of the analytic lift.
In finite dimension that range is the orthogonal complement of its
kernel. In reduced coordinates the complex lift sends a basis vector
to the analytic lift of the corresponding rational form. Its matrix
$T_j$ is the matrix of these images in the chosen target basis.
Rationality of its entries and the common image condition give
$w_{j,0}(T_j z)=\theta'_j(j_L z)$ in these coordinates, with
$j_L$ including the reduced source. If $T_j z=0$, the included vector
lies in the kernel and its orthogonal complement, so it vanishes;
the target basis spans the prescribed image. This gives the
bijection on the reduced spaces.

Three normalizations enter here. The Petersson pairing on multiplicity
maps is the integral of their matrix coefficients at unit representation
vectors over the chosen fundamental domain. The archimedean frame
multiplies this pairing by $\prod_\rho\dim K_\rho$, where $K_\rho$
is the lowest $K$-type at $\rho$. For a finite level representation
with group $H$ the adelic adjoint is $|H|^{-1}T^*$; at the trivial
level this agrees with the ordinary adjoint. In the scalar formula
the local index $I$ satisfies $\operatorname{vol}([V])=I\,\mathrm{Mass}(V)$
and $I\|f\|^2=\prod_{v\in S}d_v^{-1}$ with the specified test vector
and measures. Substitution cancels $I$ from this product, giving
$\operatorname{vol}([V])\|f\|^2/(\prod_{v\notin S}d_v)
=\mathrm{Mass}(V)/(\prod_v d_v)$.

At another embedding $\sigma$ representing $v$, the source uses a
lower bound on the reduced quadratic form. Write $s>0$ for the size of
the finite quotient of the chosen level, $C_\sigma$ for the sum of
the conjugate source pairings, and $J_{W,\sigma}$ for the realization
of the target coordinates. For a rational reduced vector $z$, the two
forms in that estimate are
\begin{align*}
 h_{T,j,v}(z,z)
   &=s^{-1}\operatorname{Re} C_\sigma(\mathrm{toF}\,z,\mathrm{toF}\,z),\\
 h_{W,j,v}(T_jz,T_jz)
   &=s^{-1}\|J_{W,\sigma}(\operatorname{coord}(T_jz))\|^2 .
\end{align*}
In the second line the realization scales the basis vectors by
$s^{-1/2}$. In the first line the majorant on the rational reduced
space is expressed in that same quotient normalization. The conjugate
lift estimate therefore gives $h_{W,j,v}(T_jz,T_jz)\ge
\exp(b_j)h_{T,j,v}(z,z)$, where
\[
 b_j=\mathrm{hookPlus}_j-8fn_V^2 L_p(B),\qquad
 L_p(B)=\log\!\bigl(B\,\operatorname{rd}_E\,g_0g_1 n_V\bigr).
\]
Explicitly, if $\lambda^+_{j\rho}$ is the Young diagram obtained by
sorting the $p_{j\rho}$ positive-side row lengths of the record, put
$H^+_{j\rho}=\prod_{\square\in\lambda^+_{j\rho}}h(\square)$, where
$h(\square)$ is one plus the number of boxes to its right in the same row
and below it in the same column.
Then $\mathrm{hookPlus}_j=\sum_\rho\log H^+_{j\rho}$.
At the base embedding the analytic main term contains the product
$H^+_{j\rho}H^-_{j\rho}$ for the two signature parts. The negative-side
hook product is a positive integer. Dropping its logarithm gives the
one-sided bound used here, including for an empty diagram.
For an active rational part, insert one nonzero vector into the scalar
identity in these forms and divide by its positive source norm.
This gives $\log\lambda_j(p,v)\ge b_j$. The selected realization uses
the lower bound at every off-base place and the positive eigenvalue
estimate at the base. The row scalar for the same $p$ uses precisely
this $\lambda_j$ in its product; both scalar identities belong to
one record. Thus transferring the scalar between the two spaces in
that record only requires the shared label and the rank equality.

For the rows (P1)--(P4) and the grid of Proposition~\ref{r6:R6:prop:nodes}, take $B=b$. The nodes $e_j=w_{j\rho}$ are positive and at most $2b$ for large $x$. The small rows satisfy $0\le v_l\le4k\Delta\le b_0/x\le2b/x\le b$. Within one embedding the grid gives $|b_l-e_j|\le mh_0\le b$, and (D5) gives $(kb)^8\le g_i$. These verify the size conditions for the indicated row parameters. The ranks give
$n_V=k+m-1+t_r\le N$, since $N-n_V=(m-1)t_r+1$.
Lemma~\ref{r5:lem:primesizes}(i),(iv) gives $N^2\le b$ and
$\log(g_0g_1)\le18L_\varepsilon$ at sufficiently large scale. Then
$\log n_V\le L_\varepsilon/2$ and
\[
 L_p(b)\le\tfrac{39}{2}L_\varepsilon+\log\operatorname{rd}_E
     \le20\Lambda_a .
\]
Here $L_\varepsilon=\log b$ and $\Lambda_a=\log(xb\operatorname{rd}_E)$.
Consequently the error summed over the lifts is bounded by
$160mfN^2\Lambda_a$. To apply this estimate to the earlier period
chain one must identify its row inclusion, type decomposition and
level forms with the selected record, and compare $\mathrm{hookPlus}_j$ with $\sum_\rho h_{j\rho}$.
For example, if the height uses another source frame, equality of its
Gram matrix in the specified rational basis with $X_{T,j}$ is needed
before applying the displayed quadratic inequality to that frame.
